%% 04_evaluacion.tex — Resultados de la barrida de ablación.
%%
%% Contenido numérico: testbed/experiment_results.csv (graduado por el
%% oráculo) + testbed/analysis/{summary_by_run.csv,stats.json} + figuras de
%% paper/figures/, todo producido por tools/analyze_results.py.
%% Las cifras en prosa provienen de sections/_datos.tex.

\section{Evaluación y resultados}
\label{sec:resultados}

La barrida produjo \NumCorridas{}~ejecuciones calificadas: los tres
presupuestos $\{20, 50, \infty\}$ con las cuatro configuraciones
(\NumBloquesCompletos{}~bloques completos) más el presupuesto de
\num{10}~peticiones con \code{baseline} y \code{no-interleave}. El presupuesto
de \num{10} para \code{no-priority} y \code{no-runtime-confirm} no llegó a
calificarse (una de las dos trazas quedó truncada), por lo que ese bloque se
trata como parcial y se excluye de las pruebas de significancia
(\cref{sec:res-stats}). Todos los valores proceden de
\code{tools/analyze\_results.py}.

\subsection{Eficiencia: detección frente a presupuesto}
\label{sec:res-eficiencia}

La \cref{fig:deteccion-vs-presupuesto} muestra que la exhaustividad del
\code{baseline} crece de \num{0.216} con \num{10}~peticiones a su valor de
saturación \ValorRecallTecho{} ya con \num{20}~peticiones por parámetro, y no
mejora al ampliar el presupuesto a \num{50} ni al eliminarlo por completo. Las
curvas de \code{no-interleave} y \code{no-runtime-confirm} son
indistinguibles de la del \code{baseline}; solo \code{no-priority} se separa,
y únicamente en el presupuesto intermedio de \num{20}
(\cref{fig:impacto-ablacion}).

\begin{figure}[t]
  \centering
  \IfFileExists{figures/fig1_detection_vs_budget.pdf}{%
    \includegraphics[width=\columnwidth]{fig1_detection_vs_budget.pdf}%
  }{%
    \fbox{\parbox[c][3.2cm][c]{0.9\columnwidth}{\centering\footnotesize
      \textsc{Marcador de posición} ---
      \code{figures/fig1\_detection\_vs\_budget.pdf}}}%
  }
  \caption{Exhaustividad (instancias vulnerables únicas detectadas $/$ total)
  en función del presupuesto de peticiones \code{-{}-max-payloads}
  ($\infty$ = sin límite), una curva por configuración del \emph{pipeline}.
  Las curvas de \code{baseline}, \code{no-interleave} y
  \code{no-runtime-confirm} se solapan.}
  \label{fig:deteccion-vs-presupuesto}
\end{figure}

\begin{figure}[t]
  \centering
  \IfFileExists{figures/fig2_ablation_impact.pdf}{%
    \includegraphics[width=\columnwidth]{fig2_ablation_impact.pdf}%
  }{%
    \fbox{\parbox[c][3.2cm][c]{0.9\columnwidth}{\centering\footnotesize
      \textsc{Marcador de posición} ---
      \code{figures/fig2\_ablation\_impact.pdf}}}%
  }
  \caption{Caída de exhaustividad ($\text{baseline} - \text{ablación}$) al
  desactivar cada optimización, por presupuesto. La $\Delta$ FPR es
  idénticamente nula en todas las celdas y se omite de la gráfica. El único
  efecto apreciable es \code{no-priority} en el presupuesto de \num{20}.}
  \label{fig:impacto-ablacion}
\end{figure}

\subsection{Resumen de métricas y AUC}
\label{sec:res-auc}

Al presupuesto de \num{20}~peticiones por parámetro, el \code{baseline}
alcanza un AUC de la curva de detección de \ValorAUCBaseline{}
(\cref{eq:auc}), con precisión \ValorPrecisionBaseline{}, exhaustividad
\ValorRecallBaseline{} y $F_1$ \ValorFOneBaseline{}. Desactivar el entrelazado
por contexto o la confirmación en tiempo de ejecución no cambia ninguna de
estas métricas (AUC \ValorAUCNoInterleave{} y \ValorAUCNoRuntimeConfirm{}
respectivamente). Desactivar la priorización sí las degrada: el AUC baja a
\ValorAUCNoPriority{}, la exhaustividad cae \DeltaRecallNoPriority{} y la
precisión \DeltaPrecisionNoPriority{}, lo que en términos absolutos equivale a
\NoPriorityVulnPerdidas{}~vulnerabilidades no detectadas y
\NoPriorityFPnuevos{}~falsos positivos adicionales frente al \code{baseline}.
En los presupuestos de \num{50} y $\infty$ el efecto se disuelve: las cuatro
configuraciones convergen a exhaustividad \ValorRecallTecho{} y difieren en
precisión en menos de \num{0.3}~p.p.\ (\cref{tab:auc}). La tasa de falsos
positivos es \ValorFPRBaseline{} (\num{2} de \NumTrampas{}~trampas) en todas
las configuraciones y presupuestos $\ge \num{20}$.

\begin{table}[t]
  \centering
  \caption{Métricas por ejecución y AUC de la curva de detección
  (\cref{eq:auc}). Fuente: \code{summary\_by\_run.csv}.}
  \label{tab:auc}
  \footnotesize
  \setlength{\tabcolsep}{4pt}
  \begin{tabular}{ll S[table-format=1.3] S[table-format=1.3] S[table-format=1.3] S[table-format=1.3]}
    \toprule
    Presup. & Config. & {Prec.} & {Rec.} & {$F_1$} & {AUC} \\
    \midrule
    \multirow{2}{*}{10\textsuperscript{$\dagger$}}
      & baseline           & 0.667 & 0.216 & 0.327 & 0.119 \\
      & no-interleave      & 0.667 & 0.216 & 0.327 & 0.119 \\
    \midrule
    \multirow{4}{*}{20}
      & baseline           & 0.619 & 0.644 & 0.631 & 0.338 \\
      & no-interleave      & 0.619 & 0.644 & 0.631 & 0.338 \\
      & no-priority        & 0.561 & 0.613 & 0.586 & 0.331 \\
      & no-runtime-confirm & 0.619 & 0.644 & 0.631 & 0.338 \\
    \midrule
    \multirow{4}{*}{50}
      & baseline           & 0.573 & 0.644 & 0.607 & 0.339 \\
      & no-interleave      & 0.573 & 0.644 & 0.607 & 0.340 \\
      & no-priority        & 0.571 & 0.644 & 0.605 & 0.341 \\
      & no-runtime-confirm & 0.571 & 0.644 & 0.605 & 0.339 \\
    \midrule
    \multirow{4}{*}{$\infty$}
      & baseline           & 0.573 & 0.644 & 0.607 & 0.341 \\
      & no-interleave      & 0.573 & 0.644 & 0.607 & 0.341 \\
      & no-priority        & 0.573 & 0.644 & 0.607 & 0.341 \\
      & no-runtime-confirm & 0.573 & 0.644 & 0.607 & 0.341 \\
    \bottomrule
  \end{tabular}
  \\[2pt]
  \footnotesize\textsuperscript{$\dagger$}\,Bloque parcial: \code{no-priority}
  y \code{no-runtime-confirm} no se calificaron a este presupuesto.
\end{table}

\subsection{Significancia estadística}
\label{sec:res-stats}

Sobre los vectores de coste de detección por instancia (bloque agrupado de los
\NumBloquesCompletos{}~presupuestos completos, $n = \ValorFriedmanN{}$), la
prueba de Friedman rechaza la equivalencia de las cuatro configuraciones
($\chi^2 = \ValorFriedmanChiSq{}$, $p = \ValorFriedmanP{}$). Sin embargo, el
tamaño del efecto es \emph{insignificante} en las tres comparaciones pareadas
\code{baseline} frente a ablación: $\delta$ de Cliff de
\ValorCliffNoInterleave{} (\code{no-interleave}), \ValorCliffNoPriority{}
(\code{no-priority}) y \ValorCliffNoRuntimeConfirm{}
(\code{no-runtime-confirm}), todos con $\lvert\delta\rvert < \num{0.02}$
(\cref{tab:stats}). El post-hoc de Nemenyi no logra separar
\code{no-runtime-confirm} del \code{baseline}
($p = \ValorNemenyiNoRuntimeConfirm{}$). La lectura es que las diferencias que
Friedman detecta son de \emph{ordenación fina} de las peticiones ---
estadísticamente reales por el gran número de bloques, pero sin traducción
práctica en el conjunto finalmente detectado. El único cambio con
consecuencias operativas es la pérdida de selectividad de \code{no-priority}
en presupuestos ajustados, ya descrita en la \cref{sec:res-auc}; nótese que
\code{no-priority} incluso reduce el coste medio de detección
(\ValorCosteBaseline{} $\to$ \ValorCosteNoPriority{} peticiones), es decir,
acierta antes en promedio pero a costa de la precisión.

\begin{table}[t]
  \centering
  \caption{Pruebas sobre el coste de detección por instancia, \code{baseline}
  frente a cada ablación (bloque agrupado, $n = \ValorFriedmanN{}$). Wilcoxon
  con corrección de Benjamini--Hochberg; $p_{\text{Nem.}}$ del post-hoc de
  Nemenyi; $\delta$ de Cliff como tamaño del efecto. Fuente:
  \code{stats.json}.}
  \label{tab:stats}
  \footnotesize
  \setlength{\tabcolsep}{4pt}
  \resizebox{\columnwidth}{!}{%
  \begin{tabular}{lccccl}
    \toprule
    \code{baseline} vs & Coste medio & $p_{\text{BH}}$ & $p_{\text{Nem.}}$
      & $\delta$ Cliff & Magnitud \\
    \midrule
    no-interleave      & \num{23006}$\to$\num{23087} & \ValorWilcoxonPNoInterleave{}
      & \ValorNemenyiNoInterleave{} & \ValorCliffNoInterleave{}
      & \MagnitudCliffNoInterleave{} \\
    no-priority        & \num{23006}$\to$\num{22818} & \ValorWilcoxonPNoPriority{}
      & \ValorNemenyiNoPriority{} & \ValorCliffNoPriority{}
      & \MagnitudCliffNoPriority{} \\
    no-runtime-confirm & \num{23006}$\to$\num{22997} & \ValorWilcoxonPNoRuntimeConfirm{}
      & \ValorNemenyiNoRuntimeConfirm{} & \ValorCliffNoRuntimeConfirm{}
      & \MagnitudCliffNoRuntimeConfirm{} \\
    \bottomrule
  \end{tabular}}
  \\[2pt]
  \footnotesize\emph{Friedman global: $\chi^2 = \ValorFriedmanChiSq{}$,
  $p = \ValorFriedmanP{}$.}
\end{table}

\subsection{Diagrama de diferencias críticas}
\label{sec:res-cd}

La \cref{fig:cd} ordena las cuatro configuraciones por su rango medio de coste
de detección (bloque agrupado, $n = \ValorFriedmanN{}$). El post-hoc de
Nemenyi solo deja sin separar a \code{baseline} y \code{no-runtime-confirm}
(barra horizontal); \code{no-interleave} queda en el peor rango
(detección más tardía) y \code{no-priority} en el mejor. Este último punto
debe leerse con cuidado: \code{no-priority} adelanta la primera detección en
promedio porque prueba antes \emph{payloads} triviales de alto rendimiento
aparente, pero ---como muestra la \cref{tab:auc}--- lo hace a costa de
\NoPriorityVulnPerdidas{}~vulnerabilidades y \NoPriorityFPnuevos{}~falsos
positivos. El diagrama compara \emph{cuándo} detecta cada configuración, no
\emph{cuánto} ni con qué precisión.

\begin{figure}[t]
  \centering
  \IfFileExists{figures/fig3_cd_diagram.pdf}{%
    \includegraphics[width=\columnwidth]{fig3_cd_diagram.pdf}%
  }{%
    \fbox{\parbox[c][2.6cm][c]{0.9\columnwidth}{\centering\footnotesize
      \textsc{Pendiente} --- \code{figures/fig3\_cd\_diagram.pdf}.}}%
  }
  \caption{Diagrama de diferencias críticas (post-hoc de Nemenyi,
  $\alpha = 0{,}05$) sobre el rango medio de coste de detección por instancia.
  Menor rango $=$ detección más temprana. Las configuraciones unidas por una
  barra horizontal no son distinguibles. Dado que todos los tamaños del efecto
  son insignificantes (\cref{tab:stats}), el diagrama es indicativo del orden,
  no de una ventaja práctica.}
  \label{fig:cd}
\end{figure}

\subsection{Alcance: motor determinista y agente de triaje}
\label{sec:res-alcance}

La barrida de ablación anterior se ejecutó sobre el \textbf{motor determinista}
(\code{-{}-ai-no}), de forma deliberada: aislar el efecto de la planificación de
\emph{payloads} exige que la decisión de detección no dependa de un componente
estocástico. El agente de triaje de la \cref{sec:agente} se evalúa por una vía
independiente y complementaria ---la puerta de aceptación de la
\cref{sec:ia-validacion}--- que mide su capacidad de suprimir falsos positivos
sin incurrir en falsos negativos sobre el conjunto disjunto de validación, con
la tasa de falsos negativos (FNR) como magnitud gobernante por su coste
asimétrico. Mantener ambas evaluaciones separadas evita confundir la eficiencia
del planificador con la fiabilidad del clasificador, y permite reproducir cada
contribución por su cuenta: la ablación sobre el banco contenedorizado, y la
alineación del agente sobre infraestructura de aceleración estándar con escalado
dinámico de parámetros según la memoria disponible.

\subsection{Síntesis de hallazgos}
\label{sec:res-sintesis}

\begin{enumerate}
  \item El \code{pipeline} \code{baseline} satura en exhaustividad
  \ValorRecallTecho{} con solo \num{20}~peticiones por parámetro y FPR
  \ValorFPRBaseline{}; ampliar el presupuesto no aporta detección, solo
  falsos positivos (la precisión cae de \ValorPrecisionBaseline{} a
  \num{57.3}\% entre los presupuestos \num{20} y $\infty$).
  \item La \textbf{priorización por confianza} es la única etapa cuya
  desactivación tiene coste medible, y se concentra en presupuestos
  ajustados: $-\DeltaRecallNoPriority{}$ de exhaustividad y
  $-\DeltaPrecisionNoPriority{}$ de precisión a presupuesto \num{20}, efecto
  que desaparece a partir de \num{50}.
  \item El \textbf{entrelazado por contexto} y la \textbf{confirmación en
  tiempo de ejecución} no alteran la exhaustividad, la precisión ni la FPR en
  este banco; su efecto sobre el orden de las peticiones es estadísticamente
  detectable pero de magnitud insignificante.
  \item En consecuencia, el valor de la planificación de \emph{payloads} en
  este escenario es de \emph{anticipación} (mayor AUC bajo presupuesto
  reducido), no de \emph{techo} de detección.
\end{enumerate}
